\songtitle{\textkorean{애매한 선} (Blurry Line)}

\tag*{2026}\tag{k-pop}\tag{korean-lyrics}\tag{thomas-entourage}\tag{natalia-kyou-park}\label{sec:blurry-line}
% -> entourage-eco.tex
\begin{notes}
Lyrics by Claude from a draft in English by Carlos Thompson, from the \emph{Thomas' Entourage} narrative universe, in the voice of Natalia -- things slightly out of place at home that Michael would never have touched, a strange scent on his travel clothes, a hidden folder on his iPhone she will not open, and the ambiguity she leans on too. \emph{Recalibrate the Space} (page \pageref{sec:recalibrate-the-space}) sets Grok's lyrics from the same draft. Retrieved from Gabriel's Suno page (V6-MINI, 18 November 2026).
\end{notes}
\begin{style}
K-pop, mid-tempo groove, female lead in a low, dark, nasal, reedy register with a layered chorus, clean theatrical production, polished synths, punchy drums, and resonant bass.
\end{style}
\begin{lyrics}
\songpart{Verse 1}
\begin{korean}
돌아온 집엔 자리가 달라진 걸\\
마이클이 건드릴 리가 없는 걸\\
여행에서 낯선 향기 밴 옷깃에\\
나도 떠나 공항 호텔 아닌 곳에
\end{korean}

\songpart{Pre-Chorus}
\begin{korean}
확실히는 알 수 없어\\
묻기도 알기도 싫어\\
선 얘기는 한 적 없고\\
그 얘기는 하기 싫고
\end{korean}

\songpart{Chorus}
\begin{korean}
나도 이 모호함에 기대고 있어\\
마이클을 사랑해 오해하진 마\\
하지만 나를 다시 맞추고 싶어\\
우리 중심 안 닿으면 바람 아냐
\end{korean}

\songpart{Verse 2}
\begin{korean}
그의 아이폰에 숨은 폴더 하나\\
여는 법을 알아도 난 열지 않아\\
우리 집에선 그가 내 곁에 있어\\
그거면 돼 나머진 다 상관없어
\end{korean}

\songpart{Pre-Chorus}
\begin{korean}
확실히는 알 수 없어\\
묻기도 알기도 싫어\\
말 안 하면 선도 없고\\
그 얘기는 하기 싫고
\end{korean}

\songpart{Chorus}
\begin{korean}
나도 이 모호함에 기대고 있어\\
마이클을 사랑해 오해하진 마\\
하지만 나를 다시 맞추고 싶어\\
집이 여전히 집이면 바람 아냐
\end{korean}

\songpart{Bridge}
\begin{korean}
다음 단계를 얘기하게 된다면\\
그땐 마주 앉아 이야기해야지\\
서로에게 선을 먼저 긋는다면\\
아니면 계속 열어 두기로 할지\\
이제 정말 진지해야 할 때일지\\
다음 단계로 우리 정말 간다면
\end{korean}

\songpart{Final Chorus}
\begin{korean}
난 모호함에 더는 안 기댈 거야\\
사랑해 마이클 의미가 될 거야\\
이제 우리 선에 이름 붙일 거야\\
그때부터 바람 뜻도 정할 거야
\end{korean}
\end{lyrics}

